\documentclass[11pt,a4paper]{ctexart}
\usepackage[margin=2.25cm]{geometry}
\usepackage{amsmath,amssymb,amsthm,booktabs,hyperref,microtype}
\hypersetup{colorlinks=true,linkcolor=black,urlcolor=blue,citecolor=black}
\setlength{\parindent}{2em}\setlength{\parskip}{0.2em}\linespread{1.06}
\newtheorem{theorem}{定理}\newtheorem{lemma}{引理}\renewcommand{\proofname}{证明}
\newcommand{\vthree}{v_3}\newcommand{\Fthree}{F_3}
\title{\bfseries 具有指定七个三进赋值差的\\四素数构造}
\author{研究札记 · F3-PAT-1}\date{2026 年 9 月 29 日}
\begin{document}\maketitle
\begin{abstract}
设 $\Fthree(m)=\vthree(m+1)-\vthree(m-1)$。本文证明存在无穷多个正整数参数对 $(n,d)$，使
\[
n,\ n+38d,\ n+92d,\ n+146d
\]
全为素数，其单点 $F_3$ 值为 $(1,-1,-1,-1)$，三个以 $n$ 为中心的乘积的 $F_3$ 值为 $(3,5,3)$。证明先用模 $3^6$ 的一个剩余类确定七个精确赋值，再应用 Green--Tao 关于复杂度至多 $2$ 的素数线性形式的无条件定理。两个参数可以同时任意大；固定步长的无穷性不在本文结论之内。
\end{abstract}
\noindent\textbf{关键词：} 三进赋值；素数线性形式；局部可解性；奇异乘积。

\section{定义与主定理}
对正整数 $a\ge1$，定义
\[
\vthree(a)=\max\{k\in\mathbb Z_{\ge0}:3^k\mid a\}.
\]
对整数 $m\ge2$，定义有符号差分
\[
\Fthree(m)=\vthree(m+1)-\vthree(m-1)\in\mathbb Z.
\]
这里不涉及 $\vthree(0)$；减法在整数中进行。

\begin{theorem}[F3-PAT-1]
存在无穷多个整数参数对 $(n,d)$，其中 $n\ge2$、$d\ge1$，使
\[
p_0=n,\quad p_1=n+38d,\quad p_2=n+92d,\quad p_3=n+146d
\]
全部为素数，且同时满足
\[
(\Fthree(p_0),\Fthree(p_1),\Fthree(p_2),\Fthree(p_3))=(1,-1,-1,-1),
\tag{1}
\]
\[
(\Fthree(p_0p_1),\Fthree(p_0p_2),\Fthree(p_0p_3))=(3,5,3).
\tag{2}
\]
可以额外要求
\[
n\equiv5\pmod{729},\qquad d\equiv1\pmod{729}.\tag{3}
\]
更精确地，对每个整数 $B\ge1$，都存在满足上述全部条件的 $n>B$、$d>B$。
\end{theorem}
因为 $d>0$，四个素数严格递增，因而互不相同。条件 (3) 只足以保证七个赋值等式；全素数参数的无穷性需要下文的全局定理。

\noindent\textbf{文稿状态。} 本结果在仓库中标记为 \texttt{PAPER-AUDITED}。本文是人类可读的书面证明；尚未形式化为 Lean 定理。PDF 与 Markdown 证明包不以有限计算或仓库 CI 代替主结论的 Lean 内核证明。

\section{用一个同余类确定七个赋值}
\begin{lemma}[精确赋值的同余稳定性]
设 $A,B$ 为正整数、$R\ge1$。若
\[
A\equiv B\pmod{3^R},\qquad \vthree(B)=e<R,
\]
则 $\vthree(A)=e$。
\end{lemma}
\begin{proof}
写 $B=3^eb$，其中 $3\nmid b$。存在整数 $z$ 使
\[
A=3^e\bigl(b+3^{R-e}z\bigr).
\]
因 $R-e\ge1$，括号内模 $3$ 同余于 $b$，不被 $3$ 整除。
\end{proof}

令 $Q=729=3^6$，并假设 $n\ge2,d\ge1$ 满足 (3)。令
\[
(h_0,h_1,h_2,h_3)=(0,38,92,146),\qquad p_i=n+h_id.
\]
则
\[
(p_0,p_1,p_2,p_3)\equiv(5,43,97,151)\pmod Q.
\]
代表元两侧的赋值如下：
\begin{center}
\begin{tabular}{c|ccc}
\toprule
$b$ & $\vthree(b+1)$ & $\vthree(b-1)$ & $\Fthree(b)$\\
\midrule
5&1&0&1\\43&0&1&-1\\97&0&1&-1\\151&0&1&-1\\
\bottomrule
\end{tabular}
\end{center}
这些赋值都小于 $6$。对每个 $p_i\pm1$ 应用引理 1，即得 (1)。

三个乘积加一满足
\[
p_0p_1+1\equiv216=3^3\cdot8\pmod Q,
\]
\[
p_0p_2+1\equiv486=3^5\cdot2\pmod Q,
\]
\[
p_0p_3+1\equiv756=3^3\cdot28\pmod Q.
\]
其代表元的赋值依次为 $3,5,3$，均小于 $6$，所以引理 1 给出对应邻数的精确赋值。另一方面，对 $i=1,2,3$，
\[
p_0\equiv2,\quad p_i\equiv1\pmod3
\quad\Longrightarrow\quad
p_0p_i-1\equiv1\pmod3.
\]
因此 $\vthree(p_0p_i-1)=0$，得到 (2)。这一部分对所有满足前提的整数参数成立，完全不需要素数性。

\section{四个仿射线性形式}
作代换
\[
n=5+Qu,\qquad d=1+Qv.\tag{4}
\]
定义 $L_i(u,v)=5+h_i+Qu+Qh_iv$，即
\[
\begin{aligned}
L_0(u,v)&=5+729u,\\
L_1(u,v)&=43+729u+27702v,\\
L_2(u,v)&=97+729u+67068v,\\
L_3(u,v)&=151+729u+106434v.
\end{aligned}
\]
记 $L=(L_0,L_1,L_2,L_3)$。在凸方框 $K_X=[X,2X]^2$ 内计数。四个形式系数固定，并满足
\[
729X\le L_i(u,v)\le215000X\qquad((u,v)\in K_X,X\ge1).\tag{5}
\]
故所有形式为正且 $\log L_i\sim\log X$。

\subsection{复杂度至多 2}
线性部分的系数向量为 $Q(1,h_i)$。若 $i\ne j$，则
\[
\det\begin{pmatrix}Q&Qh_i\\Q&Qh_j\end{pmatrix}=Q^2(h_j-h_i)\ne0.
\]
所以不同形式之间不存在 $L_i=aL_j+b$（$a,b\in\mathbb Q$）这样的仿射依赖。固定任意 $i$，把另外三个形式各自放入一个单元素组，则 $L_i$ 不在任一组的仿射线性张成中。依照 Green--Tao 的定义，系统复杂度至多 $3-1=2$。

\subsection{每个素数处的局部可解性}
对素数 $\ell$，定义局部因子
\[
\beta_\ell=\frac1{\ell^2}\sum_{u,v\bmod\ell}
\prod_{i=0}^3\left(\frac{\ell}{\ell-1}\mathbf 1_{\ell\nmid L_i(u,v)}\right).\tag{6}
\]
当 $\ell=3$ 时，四形式恒为 $(2,1,1,1)\pmod3$，故 $\beta_3=(3/2)^4>0$。

当 $\ell\ne3$ 时，$Q$ 在模 $\ell$ 下可逆。取
\[
u\equiv-4Q^{-1}\pmod\ell,\qquad v\equiv-Q^{-1}\pmod\ell.
\]
由 (4) 得 $n\equiv1,d\equiv0\pmod\ell$，于是四个 $L_i\equiv1\pmod\ell$。因此每个素数处都存在使四形式同时非零的参数，从而 $\beta_\ell>0$。

\section{奇异乘积与全素数计数}
\subsection{奇异乘积严格为正}
各局部因子为正之外，还须检查无穷乘积尾部。对 $\ell>146$，四个 $h_i$ 模 $\ell$ 互异，且 $(u,v)\mapsto(n,d)$ 是模 $\ell$ 的双射。当 $d=0$ 时允许的 $n$ 有 $\ell-1$ 个；每个 $d\ne0$ 排除四个不同的 $n=-h_id$，所以允许的 $n$ 有 $\ell-4$ 个。允许参数对总数为
\[
(\ell-1)+(\ell-1)(\ell-4)=(\ell-1)(\ell-3).
\]
代入 (6) 得
\[
\beta_\ell
=\frac{(\ell-1)(\ell-3)}{\ell^2}\left(\frac\ell{\ell-1}\right)^4
=\frac{\ell^2(\ell-3)}{(\ell-1)^3}
=1-\frac{3\ell-1}{(\ell-1)^3}
=1+O(\ell^{-2}).\tag{7}
\]
于是 $\sum_{\ell>146}|\beta_\ell-1|<\infty$，尾部乘积收敛且非零。结合有限个小素数处的正性，得到
\[
\mathfrak S_L:=\prod_{\ell\ \mathrm{prime}}\beta_\ell\in(0,\infty).\tag{8}
\]

\subsection{应用 Green--Tao 的无条件推论}
本文唯一的全局外部输入是 Green--Tao 2010 的 Corollary 1.7：复杂度至多 $2$ 的仿射线性形式系统满足其素数计数渐近式。这里使用的推论本身是无条件的。

四个形式在 $K_X$ 上均为正，故实数处因子为
\[
\beta_\infty=\operatorname{vol}_2(K_X)=X^2.
\]
记 $\mathbb P$ 为素数集合，并令
\[
A(X)=\#\{(u,v)\in K_X\cap\mathbb Z^2:L_i(u,v)\in\mathbb P\ (0\le i\le3)\}.
\]
结合尺寸、复杂度、局部因子检查，得到
\[
A(X)=(\mathfrak S_L+o(1))\frac{X^2}{(\log X)^4}.\tag{9}
\]
这里确实数素数而非所有素数幂；由 (8)，$A(X)\to\infty$。

\subsection{完成主定理}
将每个被 $A(X)$ 计数的参数对映为
\[
(n,d)=(5+729u,1+729v).
\]
该映射为单射，四个 $p_i=L_i(u,v)$ 全为素数，第 2 节保证七个赋值等式。给定任意 $B\ge1$，取充分大的 $X$ 使 $A(X)>0$ 且 $1+729X>B$，就得到 $n>B,d>B$。定理 1 及其同余加强版得证。

\section{固定步长与 Dickson 猜想}
若固定任意正整数 $d=D$，四个形式成为单变量形式
\[
n,\quad n+38D,\quad n+92D,\quad n+146D.
\]
它们两两仿射相关，属于上述复杂度定义中的无限复杂度情形，因此不能使用复杂度至多 $2$ 的推论。

特别地，固定 $d=1$ 并保留 $n=5+729u$ 后，四形式为
\[
5+729u,\quad43+729u,\quad97+729u,\quad151+729u.
\]
它们仍没有局部障碍，因此 Dickson 猜想预言这一固定步长版本也有无穷多解，而第 2 节自动保证七个赋值；这里并未证明该预言。

本文证明的是无穷多个参数对 $(n,d)$，并且两个参数可同时任意大。这不推出存在某个固定 $D$ 使 $(n,D)$ 有无穷多解。

\section*{文稿来源与复现}
本文根据仓库已审计证明包整理，结果编号为 F3-PAT-1，状态为 \texttt{PAPER-AUDITED}。证明包包含精确陈述、完整证明、引理依赖与审计记录、形式化映射及阻塞项。有限复核脚本只用于检查算术一致性，不承担无限范围证明责任。

\section*{参考文献}
\begin{thebibliography}{9}
\bibitem{GT10} Ben Green and Terence Tao. \emph{Linear equations in primes}. Annals of Mathematics, 171 (2010), 1753--1850. DOI: 10.4007/annals.2010.171.1753. Definition 1.5、Corollary 1.7 及素数计数式见固定预印本 arXiv:math/0606088v2。
\end{thebibliography}
\end{document}
